% Generated from locale/tr/content/proof-theory/proof-search/introduction.tex; do not hand-edit.


\olfileid[tr]{pt}{ps}{int}

\olsection{Giriş}

$\Log{G3c}$ ya da $\Log{N1i}$ gibi !!{proof} sistemlerinin aksiyomları ve
çıkarım kuralları, hangi sekant ya da !!{formula} ağaçlarının doğru
!!{proof}s sayılacağını tanımlar. Böyle bir sekant ya da !!{formula} ağacı
(her adımda hangi kuralların uygulandığı veya varsayımlarla onları kapatan
çıkarım kurallarının etiketleri gibi ek bilgilerle birlikte) verildiğinde,
aksiyomlarla kuralların tanımları ağacın doğru !!{proof} olup olmadığını
denetlememizi sağlar. Verilen varsayımlardan, verilmiş bir sekant ya da
!!{formula} için !!a{proof} \emph{bulmak} ise farklı ve genellikle daha zor
bir sorundur. Bu, \emph{ispat arama} sorunudur.

İspat aramanın son derece basit bir yordamı vardır. !!{formula}{}leri sonlu bir
alfabenin simge dizileri olarak düşünebiliriz. Sekantlarla, sekant ya da
!!{formula} ağaçları da (ağacın dallanmasını göstermek için belki özel
parantezler kullanılarak) !!{formula} dizileri olarak düşünülebilir.
Aksiyomlarla kurallar, bir simgeler dizisinin doğru !!{proof} gösterip
göstermediğini mekanik olarak denetlememizi sağlar. Böylece \emph{bütün
olanaklı} doğru !!{proof}{}ları mekanik olarak üretebiliriz (!!{proof}{}ları temsil
etmekte kullanılan alfabenin bütün simge dizilerini üretip her aday dizinin
doğru bir !!{proof} gösterip göstermediğini denetleyerek). Basit ispat arama
yordamı şudur: Verilen sekantın (ya da açık varsayımları verilenlerin
arasında bulunan verilmiş !!{formula}{}ün) bir !!{proof}{}ını bulana kadar bütün
doğru !!{proof}{}ları üretin.

Daha iyi bir yöntem, elle !!a{proof} bulmak için kullanacağınız saf yöntemi
kullanmaktır: Son sekantla başlar, !!a{formula} seçer ve bir çıkarım kuralını
``geriye doğru'' uygulayarak, !!{proof}{}ları bulunsaydı son çıkarımla birleşip
verilen sekantın bir !!a{proof}{}ını oluşturacak bir ya da daha çok öncül
üretirsiniz. Benzer (fakat daha karmaşık) bir yöntem doğal tümdengelim
sisteminde !!a{proof} bulmak için kullanılabilir.

Fakat bu yöntem şaşmaz değildir: Yanlış kuralı ya da !!{formula}{}leri yanlış
sırada seçerseniz çıkmaza girebilirsiniz. Ayrıca yöntem tam olarak
belirlenmemiştir: Her aşamada seçilebilecek birden çok !!{formula} ya da
geriye doğru uygulanabilecek birden çok kural bulunabilir. \emph{İspat arama
algoritması}, varsa bir !!a{proof} bulacak (yani şaşmaz) tam belirlenmiş bir
yordamdır.

Bazı !!{proof} sistemleri ötekilerden daha basit ispat arama algoritmalarına
elverişlidir. Sekant sistemlerinde !!a{formula}{}ün !!{main operator}{}ü ile
formülün solda mı sağda mı bulunduğu, geriye doğru hangi kuralın uygulanacağını
belirler. Örneğin $!A \land !B$'nin sağda bulunduğu
$\Gamma \Sequent \Delta, !A \land !B$ sekantının bir !!a{proof}{}ını bulmak
istiyorsanız \RightR{\land} kuralını geriye doğru uygulayıp karşılık gelen
$\Gamma \Sequent \Delta, !A$ ile $\Gamma \Sequent \Delta, !B$ öncüllerini
üretmelisiniz. Ancak sisteminiz $\Log{G1c}$ ise büzülme size daha çok seçenek
verdiği için işi zorlaştırır: \RightR{\Contraction} kuralını geriye doğru
uygulayıp $\Gamma \Sequent \Delta, !A \land !B, !A \land !B$ öncülünü de
üretebilirsiniz. Ardından
$\Gamma \Sequent \Delta, !A \land !B, !A \land !B, !A \land !B$ öncülünü
ve böyle devam eden öncülleri üretirsiniz. $\Log{G2c}$'de durum daha da
kötüdür. \RightR{\land} kuralını geriye doğru uygulamak, $\Gamma =
\Gamma_1 \cup \Gamma_2$ olacak $\Gamma_1,\Gamma_2$ ile $\Delta =
\Delta_1 \cup \Delta_2$ olacak $\Delta_1,\Delta_2$'yi tahmin edip
$\Gamma_1 \Sequent \Delta_1, !A$ ve $\Gamma_2 \Sequent \Delta_2, !B$
öncülleriyle şansınızı denemenizi de gerektirir. Yanlış tahmin daha sonra
çıkmaza götürebilir; başa dönüp başka $\Gamma_1,\Gamma_2,\Delta_1,\Delta_2$
seçmeniz gerekir. !!{formula}, $!A \lor !B$ ise \RightR{\lor} için iki
olasılıktan birini seçerek $\Gamma \Sequent \Delta, !A$ ya da
$\Gamma \Sequent \Delta, !B$ öncüllerinden birini üretmelisiniz. Yanlış
seçim yine geri dönmeyi gerektirir.

$\Log{G3c}$ sisteminde bu sorunlar yoktur. Yapısal kuralları bulunmaz ve
çıkarım kuralı seçimini !!{main operator} ile !!{formula}{}ün bulunduğu taraf
belirler. Dolayısıyla $\Log{G3c}$ ispat aramaya $\Log{G1c}$ ya da
$\Log{G2c}$'den daha uygundur. Başarılı arama $\Log{G3c}$ içinde !!a{proof}
verir; bu !!{proof} daha sonra $\Log{G1c}$ ya da $\Log{G2c}$'ye (hatta
$\Log{N1c}$'ye) çevrilebilir. Dolayısıyla $\Log{G3c}$ için bir ispat arama
algoritmasıyla $\Log{G3c}$ !!{proof}{}larını öteki sistemlerdeki !!{proof}{}lara
çeviren algoritmalar, bu sistemlerin ispat arama sorununu da çözer.

Bir ispat arama algoritmasının şaşmaz olduğunu nasıl biliriz? Bunun için,
algoritma !!a{proof} bulamazsa hiçbir !!{proof}{}ın var olamayacağını
göstermeliyiz. Sonuçta kötü bir algoritma seçersek ispat var olduğu halde onu
bulamadığı durumlarla karşılaşabiliriz. Bütün olanaklı !!{proof}{}ları taradığı
için basit algoritmanın böyle bir sorunu yoktur. Fakat bir ispat arama
algoritmasının verimli olması ve olabildiğince az çaba harcaması beklenir.

İspat arama algoritmasının şaşmaz olduğunu göstermenin temel fikri şudur:
Algoritma başarısızsa, başarısızlık biçiminden sekantın hiçbir !!a{proof}{}ının
var olamayacağını göstermeliyiz. Bunu da sekantın geçerli olmadığını göstererek
yaparız. Klasik birinci dereceden mantıkta geçerli sekantlar
$\Gamma \Sequent \Delta$, her !!{structure}{}nın $\Gamma$ içindeki en az bir
!!{formula}{}yı yanlış ya da $\Delta$ içindeki en az bir !!{formula}{}yı doğru
yaptığı sekantlardır. $\Log{G3c}$ \emph{sağlam}dır; yani yalnızca geçerli
sekantları !!{prove}{}r. Bu, !!{proof}s üzerinde tümevarımla gösterilir:
Aksiyomlar açıkça geçerlidir; bir kuralın öncülleri geçerliyse sonucu da
geçerlidir. Bir ispat arama algoritmasının şaşmaz olduğunu göstermek için
$\Gamma \Sequent \Delta$ sekantının bir !!a{proof}{}ını arama başarısız olduğunda,
$\Gamma$ içindeki bütün !!{formula}{}leri doğru ve $\Delta$ içindeki bütün
!!{formula}{}leri yanlış yapan !!a{structure} tanımlayabildiğimizi gösteririz.
Bu aynı zamanda $\Log{G3c}$'nin \emph{tam} olduğunu, yani
$\Gamma \Sequent \Delta$ geçerliyse onun bir $\Log{G3c}$-!!{proof}{}ı
bulunduğunu gösterir. Başarısız bütün ispat aramaları sekantın geçersiz
olduğunu gösterdiği için geçerli bir sekanta yönelik her ispat araması
başarılı olmak zorundadır.

